\section{Firms and market structure}

\begin{qfmethod}[Cost and output decisions: breakeven and shutdown prices, the short- and long-run supply of a competitive firm, and how economies and diseconomies of scale shape costs under perfect and imperfect competition]\label{meth:cost-output-shutdown}
\qffield{Statement} \emph{Short run.} A firm that takes the price $p$ as given produces where price equals marginal cost on the rising part of the marginal cost curve \citep[pp.~416--418]{varian-microeconomics-2014}. Producing nothing still costs the fixed cost $F$, so it pays to produce only if revenue covers variable cost. The \emph{shutdown condition} is
\[ p<\mathrm{AVC}(y)=c_v(y)/y \]
at the best output: below the minimum of average variable cost, the firm loses less by closing (p.~419). Between minimum AVC and minimum average total cost the firm operates at a loss, because each sale contributes to fixed costs. At the \emph{breakeven} price, where marginal cost crosses average total cost, economic profit is zero; above it the firm earns a profit \citep[sec.~8.2]{openstax-microeconomics-3e}.

\emph{Long run.} All costs are variable, so the firm stays only if $p\geq\mathrm{AC}(y)$ (Varian, p.~426). Entry and exit drive a competitive industry to zero economic profit, with price at the minimum of long-run average cost (pp.~433--437).

\emph{Scale.} The long-run average cost curve allows every input to change. It falls with output under \emph{economies of scale}, is flat under constant returns and rises under \emph{diseconomies} (OpenStax, sec.~7.5; Varian, pp.~359--361, 405--409). The output at which long-run average cost reaches its minimum is the \emph{minimum efficient scale}. How large that scale is relative to market demand shapes the structure:
\begin{itemize}
\item when demand at that cost supports many plants of efficient size, competition can drive price to the minimum of long-run average cost;
\item when average cost still falls at the output the whole market demands, one firm can supply it most cheaply: a \emph{natural monopoly} (Varian, pp.~469--471);
\item under monopolistic competition the long-run tangency of demand with average cost lies on the falling part of the cost curve, so firms produce below minimum efficient scale and price exceeds marginal cost (pp.~494--497).
\end{itemize}

\qffield{Practice} \emph{Short run} (original). Fixed cost 108, variable cost $c_v(y)=y^3-12y^2+60y$, so $\mathrm{AVC}=y^2-12y+60$ and $\mathrm{MC}=3y^2-24y+60$:
\begin{itemize}
\item AVC is lowest at $y=6$, where it equals 24 and MC: the shutdown price;
\item average total cost is lowest at $y=7.08$, where MC equals it at 40.42: the breakeven price.
\end{itemize}
\begin{center}\small
\begin{tabular}{@{}rrrl@{}}
\toprule
Price & Output (MC $=p$) & Profit & Decision \\
\midrule
60 & 8.00 & 148.00 & produce, profit \\
40 & 7.06 & $-2.97$ & produce, loss below the 108 of closing \\
30 & 6.45 & $-70.61$ & produce, loss below the 108 of closing \\
20 & -- & $-108.00$ & shut down; producing would lose 131.29 \\
\bottomrule
\end{tabular}
\end{center}

\emph{Long run and scale} (original). Long-run average cost $2{,}000/y+20+0.004y$ falls to its minimum of 25.66 at $y=707$, the minimum efficient scale.
\begin{itemize}
\item If 4,000 units are demanded at about that cost, the market holds about 5.7 plants of efficient size, consistent with competition.
\item If only 600 are, one plant cannot even reach minimum efficient scale, and a single supplier has the lowest cost: a natural monopoly.
\item A monopolistic competitor with demand $p=33.86-0.02y$ reaches the zero-profit tangency at $y=289$ and $p=28.08$. That is 41 percent of efficient scale, at a cost 2.43 above the minimum.
\end{itemize}

\qffield{Conditions} The shutdown rule compares avoidable costs. A cost that would be paid anyway, such as a lease that cannot be broken, is fixed for this decision even if accountants classify it as variable, and one that closure avoids is variable. Breakeven uses economic cost, including the opportunity cost of the owners' capital.

\qffield{Pitfalls} Closing as soon as price falls below average total cost. Treating sunk costs as relevant to the decision to continue. Assuming a falling average cost curve means increasing returns at every output.

\qffield{See also} \hyperref[def:market-structures-pricing]{market structures}; \hyperref[meth:industry-analysis]{industry analysis}.

\end{qfmethod}

\begin{qfdefinition}[Market structures compared: perfect competition, monopolistic competition, oligopoly and monopoly, with price, output and pricing strategy under monopolistic competition and oligopoly]\label{def:market-structures-pricing}
\qffield{Statement} \emph{Characteristics} (Interpretation, assembled from \citep[chaps.~23--28]{varian-microeconomics-2014} and \citep[chap.~10]{openstax-microeconomics-3e}):
\begin{center}\small
\begin{tabular}{@{}p{0.16\linewidth}p{0.18\linewidth}p{0.18\linewidth}p{0.2\linewidth}p{0.18\linewidth}@{}}
\toprule
 & Perfect competition & Monopolistic competition & Oligopoly & Monopoly \\
\midrule
Sellers & many, small & many & few, interdependent & one \\
Product & identical & differentiated & identical or differentiated & no close substitute \\
Entry & free & easy & barriers & blocked \\
Pricing power & none: price taker & some, from differentiation & depends on rivals' responses & greatest \\
Non-price competition & none & advertising, variety & often heavy & limited \\
Long-run profit & zero & zero & can persist & can persist \\
\bottomrule
\end{tabular}
\end{center}

\emph{Monopoly and markups.} A firm facing demand elasticity $\varepsilon$ sets marginal revenue equal to marginal cost, so $p=\mathrm{MC}/(1-1/|\varepsilon|)$; it never operates where demand is inelastic \citep[pp.~458--461]{varian-microeconomics-2014}.

\emph{Monopolistic competition.} Each firm faces a downward-sloping demand curve for its own differentiated product and prices like a monopolist on it. Profits attract entry, which shifts each firm's demand down until it is tangent to average cost. In the long run price equals average cost but exceeds marginal cost, and output is below the scale that minimises average cost (pp.~494--497; OpenStax, sec.~10.1). Strategy rests on differentiation and advertising, which aim to make the firm's demand less elastic.

\emph{Oligopoly.} Each firm's best choice depends on its rivals', so outcomes depend on the strategic variable and the order of moves (Varian, pp.~516--537):
\begin{itemize}
\item \emph{Cournot}: firms choose quantities simultaneously. With $n$ firms, $p(1-s_i/|\varepsilon|)=\mathrm{MC}_i$, so a firm's markup rises with its share $s_i$, and price approaches marginal cost as $n$ grows (pp.~525--530).
\item \emph{Stackelberg}: a leader chooses its quantity first, anticipating the follower's reaction, and produces more (pp.~517--522).
\item \emph{Price leadership}: a dominant firm sets the price and smaller firms take it (pp.~522--525).
\item \emph{Bertrand}: with identical products and simultaneous price setting, price falls to marginal cost even with two firms (pp.~530--531).
\item \emph{Collusion}: a cartel maximises joint profit at the monopoly output, but each member gains by producing more, so agreements need monitoring and punishment (pp.~531--537). The incentives form a prisoner's dilemma (OpenStax, sec.~10.2).
\item \emph{Kinked demand}: if rivals match price cuts but not increases, demand is elastic above the current price and inelastic below it, so no firm has a strong incentive to change price (OpenStax, sec.~10.2).
\end{itemize}

\qffield{Practice} \emph{One market, five outcomes} (original). Demand $p=100-Y$, marginal cost 10 for every firm:
\begin{center}\small
\begin{tabular}{@{}lrrr@{}}
\toprule
Structure & Output & Price & $(p-\mathrm{MC})/p$ \\
\midrule
Monopoly or cartel & 45.0 & 55.00 & 0.82 \\
Stackelberg (leader 45, follower 22.5) & 67.5 & 32.50 & 0.69 \\
Cournot, two firms (30 each) & 60.0 & 40.00 & 0.75 \\
Cournot, five firms (15 each) & 75.0 & 25.00 & 0.60 \\
Bertrand, or perfect competition & 90.0 & 10.00 & 0.00 \\
\bottomrule
\end{tabular}
\end{center}
A cartel member that expands from 22.5 to the Cournot best response of 33.75, while the other holds 22.5, raises its own profit from 1,012.5 to 1,139.1. The other then earns 759.4, which is why cartels need enforcement.

\qffield{Conditions} The models assume known demand and costs. Which one describes a market is an empirical question about how firms actually compete: through capacity (Cournot), prices (Bertrand) or leadership.

\qffield{Pitfalls} Inferring competitive conduct from the number of firms alone; two Bertrand competitors price at cost. Treating zero long-run profit under monopolistic competition as efficiency.

\qffield{See also} \hyperref[meth:cost-output-shutdown]{cost and output}; \hyperref[meth:market-structure-identification]{identifying the structure}; \hyperref[meth:industry-analysis]{industry analysis}.

\end{qfdefinition}

\begin{qfmethod}[Identifying a firm's market structure: the evidence to use, concentration ratios and the Herfindahl--Hirschman index, their link to markups, and their limitations]\label{meth:market-structure-identification}
\qffield{Statement} \emph{Evidence.} The structure is identified from the number and size of sellers, product differentiation, barriers to entry, pricing power and the way firms respond to each other (Interpretation, from the table in \hyperref[def:market-structures-pricing]{market structures}). Concentration measures summarise the first of these:
\begin{itemize}
\item the $N$-firm \emph{concentration ratio} is the combined market share of the largest $N$ firms, usually four to eight \citep[sec.~11.1]{openstax-microeconomics-3e};
\item the \emph{Herfindahl--Hirschman index} sums the squared percentage shares of all firms, from near 0 to 10,000 for a monopoly. Squaring weights large firms more. A merger of firms with shares $a$ and $b$ raises it by $2ab$ \citep[pp.~5--6, n.~14]{merger-guidelines-2023}.
\end{itemize}
US agencies treat a market with an index above 1,800 as highly concentrated. A merger that leaves it above 1,800 and raises it by more than 100, or creates a firm with over 30 percent and an increase of more than 100, is presumed to lessen competition substantially (pp.~5--6).

\emph{Link to markups} (Derivation, from the Cournot condition \citep[pp.~529--530]{varian-microeconomics-2014}). Weighting each firm's markup $(p-\mathrm{MC}_i)/p=s_i/|\varepsilon|$ by its share gives an industry average of $\sum s_i^2/|\varepsilon|=\mathrm{HHI}/|\varepsilon|$, with HHI as a fraction. Under Cournot competition, concentration and market power move together; under Bertrand they do not.

\emph{Limitations} (OpenStax, sec.~11.1):
\begin{itemize}
\item the market must be defined first, and a narrow definition makes concentration look higher;
\item domestic shares overstate concentration when competition is global;
\item one measure assumes competitive conditions are similar across industries.
\end{itemize}
Concentration also says nothing about entry barriers, the threat of entry or how firms actually compete (Interpretation).

\qffield{Practice} \emph{An industry with eight firms} (original), shares 30, 25, 15, 10 and four of 5 percent:
\begin{itemize}
\item four-firm ratio 80 percent; HHI $900+625+225+100+4\times25=1{,}950$, highly concentrated;
\item a merger of the third and fourth firms adds $2\times15\times10=300$, to 2,250, inside the structural presumption;
\item if firms compete in quantities and industry demand elasticity is 1.5, the average markup is $0.195/1.5=13$ percent.
\end{itemize}
OpenStax's contrast shows why the index matters: five firms of 20 percent and one firm of 77 percent with 23 of 1 percent both have a four-firm ratio of 80, but HHIs of 2,000 and 5,952.

\emph{Identification.} With few interdependent sellers, a differentiated but substitutable product and costly entry, this industry is an oligopoly. Whether prices are near cost depends on whether firms compete on price or capacity, which the HHI alone cannot reveal.

\qffield{Conditions} Shares are measured in the relevant product and geographic market, by sales, capacity or units as the question requires.

\qffield{Pitfalls} Reading a high concentration ratio as proof of monopoly pricing. Comparing indexes across differently defined markets. Ignoring imports, potential entrants and buyer power.

\qffield{See also} \hyperref[def:market-structures-pricing]{market structures}; \hyperref[meth:industry-analysis]{industry analysis}; \hyperref[meth:peer-group-selection]{peer groups}.

\end{qfmethod}
